\documentclass[11pt]{article}
\usepackage[margin=1in]{geometry}
\usepackage{amsmath,amssymb,amsthm}
\usepackage{graphicx}
\usepackage{placeins}
\usepackage{booktabs}
\usepackage[expansion=false]{microtype}
\usepackage[hidelinks]{hyperref}
\setlength{\emergencystretch}{3em}

\title{\bf Fine Structure of the Wing-Transition Memory\\
on the $6N$ Skeleton: the Lag-Two Echo and\\
the $1/\log$ Decay%
\thanks{Code and measured data: \url{https://github.com/Ruqing1963/6N-wing-finestructure}}}
\author{Ruqing Chen\\
\small GUT Geoservice Inc., Montreal, Quebec, Canada\\
\small \texttt{ruqing@hotmail.com}}
\date{June 8, 2026}

\begin{document}
\maketitle

\begin{abstract}
Part~XXXIV established that the prime wing string---right ($6N{+}1$) versus left ($6N{-}1$)---is not a
first-order Markov chain but carries an effective second-order memory, the switch sequence having a
lag-one autocorrelation of $-0.033$. Two fine questions were left open, and we settle them here with the
full prime census to $10^{9}$ ($50{,}847{,}531$ steps). First: is the small lag-two echo a mere shadow
of the lag-one term, or independent? For a first-order (AR(1)) switch process the lag-two
autocorrelation would be exactly $\rho_1^2=+0.00098$; we measure $\rho_2=+0.00417$, \emph{$4.3$ times}
larger, a partial autocorrelation $\mathrm{PACF}(2)=+0.0032$ at $23\sigma$, with $\rho_{\ge3}$
vanishing. The memory therefore reaches genuinely to lag two---a positive echo riding on the negative
nearest-neighbour term---and stops. Second: how fast does the bias fade? Across thirteen logarithmic
windows the first-order deficit $\delta_1=\tfrac12-P(\text{same})$ obeys a leading $1/\log N$ law
($\delta_1\log N\approx1.05$, nearly flat; a single power fit gives an effective exponent $0.85$, the
shortfall a slowly-varying correction the range cannot resolve), and the second-order signal $|\rho_1|$
fades in lock-step, the ratio $\delta_1/|\rho_1|\approx1.7$ holding across five decades. Both are
incremental refinements of Part~XXXIV inside the Hardy--Littlewood and Oliver--Soundararajan framework;
what we add are the specific measured numbers---the lag-two excess and the decay law. This is a
measurement, not a theorem.
\end{abstract}

\section{Introduction}
Part~XXXIV \cite{PXXXIV} read the primes as a binary string of wing symbols and found the string is not
a fair coin (Cram\'er) nor a first-order Markov chain: the switch sequence
$\sigma_i=\mathbf{1}[w_{i+1}\ne w_i]$ has autocorrelation $\rho_1\approx-0.033$, an anti-persistence we
called a topological strain memory, with a weak $+0.005$ at lag two and nothing measurable beyond. Two
questions were flagged but not answered, and both are sharp, finite, and were unknown to us in advance.

\emph{(i) Is the lag-two echo real, or derived?} A symmetric first-order chain has an independent switch
sequence, but a chain that is itself merely AR(1) in the switches would still show
$\rho_k=\rho_1^{\,k}$; the lag-two value would then be the trivial $\rho_1^2\approx+0.001$. Is the
measured echo larger than that?

\emph{(ii) What is the decay law?} Oliver and Soundararajan \cite{OS} predict the bias vanishes as
$N\to\infty$; Part~XXXIV showed $P(\text{same})$ drifting from $0.404$ toward $0.5$, but fit no rate.
Does $\delta_1$ fall like $1/\log N$, like $1/\log^2 N$, like $\log\log N/\log N$? And does the
second-order signal fade by the same law or its own?

These are refinements, not a new direction: the existence of the structure is contained in the
Hardy--Littlewood $r$-tuple heuristic \cite{HL} and the Oliver--Soundararajan bias \cite{OS}, and we
prove nothing. What we contribute are the measured numbers that pin the fine structure down. We use all
$50{,}847{,}532$ primes $3<p\le10^{9}$.

\section{Setup}
Let $w_i\in\{\mathrm{R},\mathrm{L}\}$ be the wing of the $i$-th prime $>3$ and
$\sigma_i=\mathbf{1}[w_{i+1}\ne w_i]$ the switch sequence, with overall switch rate $P(\sigma)=0.5549$
to $10^{9}$. Two null models frame the questions. For the lag-two test, the null is an AR(1) switch
process, under which the autocorrelation decays geometrically, $\rho_k=\rho_1^{\,k}$; the partial
autocorrelation $\mathrm{PACF}(2)=(\rho_2-\rho_1^2)/(1-\rho_1^2)$ isolates any lag-two structure not
already implied by lag one, and is zero under the null. For the decay law we bin the steps into
quarter-decade windows of $p$ and fit $\delta_1\sim C/(\log N)^{\alpha}$.

\section{The lag-two echo is real}
The switch autocorrelations to $10^{9}$ are
\[
\rho_1=-0.03127\ (223\sigma),\quad \rho_2=+0.00417\ (30\sigma),\quad \rho_3=+0.00009\ (1\sigma),
\]
with $\rho_{\ge3}$ indistinguishable from zero (Fig.~\ref{fig:fig1}(A)). The AR(1) null predicts
$\rho_2=\rho_1^2=+0.00098$; the measured value is $4.3$ times larger. The excess is the partial
autocorrelation
\[
\mathrm{PACF}(2)=\frac{\rho_2-\rho_1^2}{1-\rho_1^2}=+0.00319\quad(23\sigma),
\]
significantly nonzero (Fig.~\ref{fig:fig1}(B)). The memory in the switch sequence therefore reaches
genuinely to the second lag---and, with $\rho_{\ge3}$ vanishing, no further. Its sign is the surprise:
the nearest-neighbour term is negative (a switch suppresses the next switch), but the lag-two term is
\emph{positive}, a faint tendency for the state two steps on to rhyme with the present---picturesquely,
the faint overcorrection echo following the strong nearest-neighbour rebound. We offer the image only
as a mnemonic for the measured signs; it is a description of two autocorrelation coefficients, not a
mechanism. This is the
step-sequence face of the third-order residual seen directly in the wing conditionals of Part~XXXIV;
here it is resolved as a signed, lag-two-localised echo. ``Effective second-order'' is thus sharpened:
the switch process is one of memory length two, dominated by the lag-one anti-persistence with a
lag-two positive echo a tenth its size.

\section{The decay law}
Binning into thirteen quarter-decade windows from $p\sim10^{5}$ to $10^{9}$
(Table~\ref{tab:decay}), the first-order deficit $\delta_1=\tfrac12-P(\text{same})$ falls from
$0.074$ to $0.053$. Three readings agree on a leading $1/\log N$ law. The product $\delta_1\log N$ is
nearly flat at $1.05$ (rising only from $1.00$ to $1.08$ across the range), so $\delta_1\approx1/\log
N$ to leading order; a single power fit $\delta_1\sim C/(\log N)^{\alpha}$ returns $\alpha=0.85$, the
shortfall from $1$ reflecting a slowly-varying (plausibly $\log\log$) correction the five-decade range
cannot resolve (Fig.~\ref{fig:fig2}(A)). The second-order signal $|\rho_1|$ fits $\alpha=1.06$ and,
crucially, fades in step with the first: the ratio $\delta_1/|\rho_1|$ holds at $1.68$ across the whole
range (Fig.~\ref{fig:fig2}(B)), drifting only from $1.6$ to $1.8$. The first- and second-order biases
are thus two amplitudes of a single fading effect, not independent decays---consistent with both being
shadows of the same Oliver--Soundararajan correction, whose leading order in the wing language is
$\sim1/\log N$.

\begin{table}[t]
\centering\small
\begin{tabular}{rccccc}
\toprule
$\log N$ & $P(\text{same})$ & $\delta_1$ & $|\rho_1|$ & $\delta_1\log N$ & $\delta_1/|\rho_1|$\\
\midrule
13.5 & 0.4261 & 0.0739 & 0.0436 & 1.000 & 1.70\\
15.8 & 0.4359 & 0.0641 & 0.0387 & 1.048 & 1.66\\
17.6 & 0.4403 & 0.0597 & 0.0350 & 1.045 & 1.71\\
19.3 & 0.4458 & 0.0542 & 0.0310 & 1.075 & 1.75\\
20.4 & 0.4470 & 0.0530 & 0.0295 & 1.082 & 1.80\\
\bottomrule
\end{tabular}
\caption{Decay of the wing bias (selected windows; full table in the repository). $\delta_1\log N$ is
nearly flat ($\sim1/\log N$ leading), and $\delta_1/|\rho_1|$ is nearly constant (first and second order
locked together).}
\label{tab:decay}
\end{table}

\begin{figure}[t]
\centering
\includegraphics[width=\textwidth]{p35_fig1.pdf}
\caption{(A) Switch-sequence autocorrelation (blue) against the AR(1) prediction $\rho_1^{\,k}$ (red).
The lag-one terms coincide by construction; the lag-two measured value exceeds the AR(1) value, and
$\rho_{\ge3}$ vanish. (B) The lag-two value decomposed: the AR(1)-derived part $\rho_1^2$, the measured
$\rho_2$, and their difference, the partial autocorrelation $\mathrm{PACF}(2)=+0.0032$ ($23\sigma$).}
\label{fig:fig1}
\end{figure}

\begin{figure}[t]
\centering
\includegraphics[width=\textwidth]{p35_fig2.pdf}
\caption{(A) The first-order deficit $\delta_1$ and the second-order $|\rho_1|$ versus $\log N$
(log--log), with a power-law guide of exponent $0.85$. (B) The two diagnostics that read the law:
$\delta_1\log N\approx1.05$ (nearly flat: leading $1/\log N$) and $\delta_1/|\rho_1|\approx1.7$ (nearly
flat: first and second order locked).}
\label{fig:fig2}
\end{figure}

\FloatBarrier
\section{Scope and honest ledger}
Two fine questions left open in Part~XXXIV now have measured answers. The lag-two echo of the switch
sequence is real, not an AR(1) shadow: at $4.3$ times the geometric prediction and $23\sigma$ in the
partial autocorrelation, with all higher lags vanishing, the switch process has memory length exactly
two---a negative nearest-neighbour term and a positive lag-two echo a tenth its size. And the bias
fades as a leading $1/\log N$ ($\delta_1\log N\approx1.05$), the second-order signal locked to the first
at a constant ratio $\approx1.7$, so the two are one fading effect seen at two orders. The ratio is not
exactly fixed---it drifts from $1.6$ to $1.8$ across the range---but its near-constancy over five
decades says the first- and second-order biases are controlled by a single fading amplitude rather than
decaying independently.

We are explicit that this is a refinement, not a discovery. The structure itself is contained in the
Hardy--Littlewood $r$-tuple heuristic \cite{HL}, and the vanishing of the bias with scale is the
Oliver--Soundararajan prediction \cite{OS}; we prove no theorem and propose no mechanism. What this note
adds is the quantitative fine structure---the size and sign of the lag-two partial autocorrelation, the
exact-two memory length, the leading $1/\log N$ decay exponent, and the locked first-to-second-order
ratio---measured against the full prime census to $10^{9}$. Everything here equals a direct count over
the primes and takes the Hardy--Littlewood heuristic as input.

\paragraph{Data and code availability.} The sieve to $10^{9}$, the switch-sequence autocorrelations and
AR(1)/partial-autocorrelation test, the full windowed decay table and fits, and the figure scripts are
openly available at \url{https://github.com/Ruqing1963/6N-wing-finestructure}.

\begin{thebibliography}{9}
\bibitem{HL} G.~H.~Hardy, J.~E.~Littlewood, \emph{Some problems of `Partitio Numerorum'; III}, Acta Math.\ \textbf{44} (1923), 1--70.
\bibitem{OS} R.~J.~Lemke Oliver, K.~Soundararajan, \emph{Unexpected biases in the distribution of consecutive primes}, Proc.\ Natl.\ Acad.\ Sci.\ USA \textbf{113} (2016), E4446--E4454.
\bibitem{PXXXIV} R.~Chen, \emph{Topological strain memory on the 6N skeleton: the effective second-order Markov dynamics of prime wing transitions} (Part XXXIV), Zenodo, doi:10.5281/zenodo.20603211 (2026).
\bibitem{review} R.~Chen, \emph{Arithmetic geodynamics on the 6N skeleton: a review of Parts I--XIX}, Zenodo, doi:10.5281/zenodo.20585301 (2026).
\end{thebibliography}

\end{document}
